\documentclass[../../../include/open-logic-section]{subfiles}

\begin{document}

\olfileid{sth}{z}{arbintersections}
\olsection{అనుబంధం: సంవృతత, ధర్మసంగ్రహం, ఛేదనం}

\olref[sth][z][milestone]{sec}లో,
\olref[sfr][][]{part}లోని
అనౌపచారిక పనిని మళ్లీ
చూసి $\Zminus$లో చేయవచ్చో
తనిఖీ చేయాలని సూచించాం.
ఆ సూచనను పాటిస్తే,
\emph{ఒక} విషయం మిమ్మల్ని
ఆపి ఉండవచ్చు: డెడెకిండ్
\emph{సంవృతాల} విశ్లేషణలో
\emph{ఛేదనం} వాడకం.

\olref[sfr][infinite][dedekind]{Closure}
నుంచి ఇది గుర్తుచేసుకోండి:
\begin{align*}
  \closureofunder{f}{o} & = \bigcap\Setabs{X}{o \in X \text{ మరియు $X$ అనేది
$f$-సంవృతం}}.
\intertext{దీని సాధారణ రూపం ఈ కింది తరహా నిర్వచనం:}
  C & = \bigcap\Setabs{X}{\phi(X)}.
\end{align*}
అయితే ఇది హెచ్చరికలా
వినిపించాలి: అమాయక
ధర్మసంగ్రహం విఫలమవుతుంది
కాబట్టి, $\Setabs{X}{\phi(X)}$
ఉంటుందని హామీ లేదు.
అందువల్ల అటువంటి
నిర్వచనాలు \emph{అన్యాయంగా}
ఉన్నట్టుగా కనిపిస్తున్నాయి.

అదృష్టవశాత్తూ, అవి
అన్యాయం కావు. లేదా
ఇప్పుడున్న రూపంలో
\emph{అన్యాయమే} అయినా,
కొంచెం కచ్చితమైన పని
చేసి వాటిని చెల్లుబాటయ్యేలా
మార్చవచ్చు. ఏ $A \neq \emptyset$కైనా
$\bigcap A$ ఎందుకు
ఉంటుందో వివరించిన
\olref[sth][z][sep]{prop:intersectionsexist}లో
ఈ పనికి సూచన ఇచ్చాం.
ఇప్పుడు దానిని
స్పష్టంగా వివరిస్తాం.

సమితుల సమానత్వ సూత్రం
ఉన్నప్పుడు,
$C$ను $\bigcap\Setabs{X}{\phi(X)}$గా
నిర్వచించాలనుకోవడం
అంటే నిజానికి ఈ కింది
షరతును సంతృప్తిపరిచే
$C$ వస్తువు కావాలనుకోవడమే:
\begin{equation}\ollabel{bicondelimarbintersection}
	\forall x(x \in C \liff \forall X(\phi(X) \lif x \in X))\tag{*}
\end{equation}
ఇప్పుడు $\phi(S)$ అయ్యే
\emph{ఏదో} సమితి $S$
ఉందనుకుందాం. అప్పుడు
\olref{bicondelimarbintersection}ను
పొందడానికి $C$ను
\emph{వేరుచేయడం} ద్వారా
ఇలా నిర్వచించవచ్చు:
\[
	C = \Setabs{x \in S}{\forall X(\phi(X) \lif x \in X)}.
\]
ఈ నిర్వచనం కోరుకున్నట్టుగా
\olref{bicondelimarbintersection}ను
ఇస్తుందని తనిఖీ చేయడాన్ని
అభ్యాసంగా వదులుతున్నాం.
%  $x$ను స్థిరపరచండి. మొదట, తిరిగి నిర్వచించినట్టు $x \in C$ అనుకుందాం; అప్పుడు $x
%  \in S$, అలాగే $\forall X(\phi(X) \lif x \in X)$; కానీ
%  $\phi(S)$ కాబట్టి, ఇది $\forall X(\phi(X)
%  \lif x \in X)$ షరతుకు తగ్గుతుంది. తిరిగి $\forall X(\phi(X) \lif
%  x \in X)$ అనుకుందాం; $\phi(S)$ కాబట్టి $x \in S$ కూడా; అందువల్ల
%  తిరిగి నిర్వచించిన ప్రకారం $x \in c$.
ఈ సాధారణ పద్ధతితో
ఛేదనాల నిర్వచనంలో
అమాయక ధర్మసంగ్రహం
వాడుతున్నట్టుగా కనిపించే
ఏ సందర్భాన్నైనా
తప్పించవచ్చు. ఈ
ఆలోచన మొదలైన ప్రత్యేక
సందర్భంలో, అంటే
$\closureofunder{f}{o}$
విషయంలో, ఇది ఇలా
పనిచేస్తుంది.
\olref[sfr][infinite][dedekind]{closureproperties}
నిరూపణను మొదలుపెడుతూ,
$o \in \ran{f}\cup\{o\}$
అనీ $\ran{f} \cup \{o\}$
$f$-సంవృతమనీ
గుర్తించాం. కాబట్టి
మనకు కావలసినదాన్ని
ఇలా నిర్వచించవచ్చు:
\[
	\closureofunder{f}{o} = \Setabs{x \in \ran{f} \cup \{o\}}{(\forall X \ni o)(X \text{ అనేది $f$-సంవృతం} \lif x \in X)}.
\]

\end{document}
